% =========================================================
% Differential Geometry — Notes Index
% =========================================================

% ---------------------------------------------------------
% Breadcrumb
% ---------------------------------------------------------
\begin{tcolorbox}[
  colback=gray!6,
  colframe=gray!40,
  arc=2pt,
  left=8pt, right=8pt, top=6pt, bottom=6pt,
  title={\small\textbf{Where You Are in the Journey}},
  fonttitle=\small\bfseries
]
\begin{center}
\small
Calculus
$\;\to\;$ Analytical Geometry
$\;\to\;$ \textbf{Differential Geometry}
$\;\to\;$ Riemannian Geometry
$\;\to\;$ Stochastic Analysis on Manifolds
$\;\to\;$ $\cdots$
\end{center}

\medskip
\noindent\textbf{How we got here.}
Calculus built the tools of local linear approximation on
$\mathbb{R}$.
Analytical geometry gave those tools a spatial stage.
Differential geometry takes the next step: it studies
geometric objects --- curves and surfaces --- not just as sets
of points, but as objects equipped with their own intrinsic
notion of calculus.

\medskip
\noindent\textbf{What this chapter builds.}
Parametric curves in $\mathbb{R}^3$; arc length, curvature,
and the Frenet--Serret frame; regular surfaces in $\mathbb{R}^3$;
the first and second fundamental forms; Gaussian and mean
curvature; geodesics; and the Gauss--Bonnet theorem as
a global structural result connecting curvature to topology.

\medskip
\noindent\textbf{Where this leads.}
The notion of a smooth manifold abstracts away the ambient
Euclidean space entirely, replacing it with coordinate charts
and transition maps.
Riemannian geometry equips manifolds with an inner product at
each tangent space, recovering notions of length, angle, and
curvature in full generality.
This is the immediate prerequisite for stochastic analysis
on manifolds.
\end{tcolorbox}
\vspace{1em}

% ---------------------------------------------------------
% Structural Role
% ---------------------------------------------------------
\begin{tcolorbox}[
  colback=gray!6,
  colframe=gray!40,
  arc=2pt,
  left=8pt, right=8pt, top=6pt, bottom=6pt,
  title={\small\textbf{Structural Role: Differential Geometry — Calculus on Curved Spaces}},
  fonttitle=\small\bfseries
]
Differential geometry is the study of geometric objects through
the lens of calculus.

The central insight is that curvature --- the failure of a space
to be flat --- can be measured \emph{intrinsically}, without
reference to any ambient space.
Gauss's Theorema Egregium is the landmark result: Gaussian
curvature is preserved under isometries and therefore belongs
to the surface itself, not to how it sits in $\mathbb{R}^3$.

The Gauss--Bonnet theorem goes further: it equates the
integral of curvature over a compact surface to a purely
topological quantity (the Euler characteristic).
This is the first glimpse of how local analytic data
constrains global topological structure, a theme that runs
through all of modern geometry and topology.
\end{tcolorbox}
\vspace{1em}

% ---------------------------------------------------------
% Structural Roadmap
% ---------------------------------------------------------
\subsection*{Structural Roadmap}

\begin{center}
\textbf{Definitions $\longrightarrow$ Main Theorems
$\longrightarrow$ Consequences and Structural Insight}
\end{center}

The global progression is:

\begin{enumerate}
  \item Parametric curves: arc length, speed, reparametrisation
  \item Curvature and torsion; the Frenet--Serret frame
  \item Regular surfaces in $\mathbb{R}^3$: coordinate charts
        and tangent planes
  \item The first fundamental form: lengths and areas on a surface
  \item The second fundamental form: normal curvature and shape operator
  \item Gaussian and mean curvature; the Theorema Egregium
  \item Geodesics: definition, existence, and examples
  \item The Gauss--Bonnet theorem
  \item Smooth manifolds: charts, atlases, and the passage to
        intrinsic geometry
\end{enumerate}

\begin{remark*}[Primary sources]
do~Carmo, \textit{Differential Geometry of Curves and Surfaces};
Lee, \textit{Introduction to Smooth Manifolds};
Pressley, \textit{Elementary Differential Geometry}.
\end{remark*}

% ---------------------------------------------------------
% Status
% ---------------------------------------------------------
\vspace{1em}
\begin{tcolorbox}[
  colback=gray!6, colframe=gray!40, arc=2pt,
  left=8pt, right=8pt, top=6pt, bottom=6pt,
  title={\small\textbf{Status: Planned}},
  fonttitle=\small\bfseries
]
\begin{center}\Large\bfseries Coming Soon\end{center}
\vspace{6pt}
\noindent Notes, proofs, and exercises will appear here in a future revision.
\end{tcolorbox}

% Future sections (uncomment as content is written):
% \input{volume-v/geometry/differential-geometry/notes/notes-curves}
% \input{volume-v/geometry/differential-geometry/notes/notes-frenet-serret}
% \input{volume-v/geometry/differential-geometry/notes/notes-surfaces}
% \input{volume-v/geometry/differential-geometry/notes/notes-first-ff}
% \input{volume-v/geometry/differential-geometry/notes/notes-second-ff}
% \input{volume-v/geometry/differential-geometry/notes/notes-curvature}
% \input{volume-v/geometry/differential-geometry/notes/notes-geodesics}
% \input{volume-v/geometry/differential-geometry/notes/notes-gauss-bonnet}
% \input{volume-v/geometry/differential-geometry/notes/notes-manifolds}
